\fontsize{14}{15}\selectfont


El uso de sistemas de software en nuestra vida diaria está aumentando exponencialmente a medida que se integran en diversos campos. Por lo tanto, los desarrolladores buscan crear sistemas de software de alta calidad para mantenerse competitivos, lo que requiere la efectividad de cada fase del ciclo de vida del desarrollo de software. Las pruebas de software son una de las fases más críticas, ya que las fallas no detectadas pueden tener repercusiones significativas, como transacciones defectuosas, interrupciones del servicio y fallas en sistemas dependientes \cite{mohsin2022span,singh2024improved}. Por lo tanto, las pruebas de software deben ser consideradas cuidadosamente para evitar fallas en entornos reales \cite{singh2024improved,ferenc_automatically_2020}.

Inicialmente, el principal desafío en la Predicción de Fallas de Software (PFS) era la falta de conjuntos de datos públicos. Sin datos, es imposible comparar los resultados de diferentes métodos. Para abordar esto, se crearon repositorios públicos como PROMISE \cite{Sayyad-Shirabad+Menzies:2005}, inspirados en el éxito del repositorio de Aprendizaje Automático de la Universidad de California en Irvine \cite{uci_ref}. Además, muchos investigadores han proporcionado sus propios conjuntos de datos para PFS \cite{li_progress_2018}. Entre los más utilizados se encuentran PROMISE, NASA, SOFTLAB, ReLink, AEEEM, ECLIPSE1 y ECLIPSE2 \cite{shepperd2013data,d2012evaluating,wu2011relink,zimmermann2007predicting,kim2011dealing}.

Estos conjuntos de datos tradicionales se centran en métricas como líneas de código y complejidad ciclomática. Aunque útiles, pueden no capturar la complejidad de los sistemas modernos. En \cite{ferenc_automatically_2020}, se propone un enfoque que recopila instantáneas de elementos de código antes y después de las correcciones de errores, utilizando 15 proyectos en Java de GitHub. Este enfoque captura métricas cuando se detecta un error y nuevamente cuando se corrige, diferenciando entre métricas de código con y sin errores.

Las investigaciones recientes indican que la calidad de los conjuntos de datos afecta significativamente el rendimiento de la predicción de errores \cite{rathi_empirical_2023}. Los problemas incluyen sesgos, ruido, alta dimensionalidad y desequilibrio de clases \cite{liu_empirical_2015,8494821}. Este documento se centra en los dos últimos aspectos para mejorar la calidad de los datos.

El problema de alta dimensionalidad, causado por características poco relevantes, se mitiga mediante la selección de características. El desequilibrio de clases, causado por la presencia desproporcionada de instancias de una clase sobre otra, se aborda mediante técnicas de muestreo. Ambos enfoques han demostrado ser útiles en estudios previos \cite{liu_empirical_2015}.

Esta investigación evalúa un enfoque de preprocesamiento de datos en dos etapas, combinando selección de características y reducción de instancias, aplicado al conjunto de datos BugHunter \cite{ferenc_automatically_2020}. El objetivo principal es determinar la mejor combinación de técnicas que permita reducir la cantidad de características y mejorar la precisión del modelo de PFS.

Se replican los resultados de \cite{ferenc_automatically_2020} usando BugHunter y el clasificador RandomForest en Weka \cite{weka}, sirviendo como línea base. Se aplican técnicas de análisis de relevancia y control de redundancia para seleccionar métricas, variando parámetros para optimizar resultados.

La investigación se enfoca en generar evidencia experimental para responder las siguientes preguntas de investigación:

\begin{itemize}
    \item \textbf{RQ1.} ¿El análisis de relevancia y control de redundancia de características mejora la identificación de errores basada en métricas de código?
    \item \textbf{RQ2.} ¿Cuáles métricas de código son buenas predictoras de errores y en qué nivel (archivo, clase, método) proporcionan mejores resultados?
    \item \textbf{RQ3.} ¿Cómo influye la variación de parámetros en los algoritmos de análisis de relevancia y control de redundancia en los resultados de los modelos de clasificación?
    \item \textbf{RQ4.} ¿Cómo influyen los algoritmos de balanceo de clases ROS, RUS y SMOTE en los resultados de los clasificadores?
    \item \textbf{RQ5.} ¿Qué diferencias existen entre el proyecto resumen y los proyectos individuales en los resultados obtenidos en las preguntas anteriores?
\end{itemize}

\textbf{Hipótesis general:} La aplicación de técnicas avanzadas de preprocesamiento de datos mejora la precisión de los modelos de predicción de fallas de software.

\textbf{Hipótesis específicas:}
\begin{itemize}
    \item \textbf{H1:} La aplicación de técnicas de selección de características sobre conjuntos de datos de métricas de software permite reducir significativamente la cantidad de características utilizadas, sin disminuir la precisión de los modelos de predicción de fallas de software.
    \item \textbf{H2:} El uso de algoritmos de balanceo de clases (ROS, RUS, SMOTE) mejora la precisión de los modelos de predicción de fallas de software en comparación con modelos entrenados en conjuntos de datos desequilibrados.
    \item \textbf{H3:} La combinación de selección de características y balanceo de clases produce modelos de predicción de fallas de software con mayor precisión que aquellos entrenados con el conjunto de datos original sin preprocesamiento.
    \item \textbf{H4:} Las métricas de código seleccionadas mediante análisis de relevancia y control de redundancia son mejores predictoras de fallas de software que el conjunto completo de métricas.
    \item \textbf{H5:} La integración de métricas de varios proyectos en un solo conjunto de datos (proyecto resumen) no disminuye significativamente la precisión de los modelos de predicción de fallas de software en comparación con el análisis individual por proyecto.
\end{itemize}

\textbf{Contribuciones del Trabajo de Investigación}

\begin{itemize}
    \item Se presenta una evaluación experimental exhaustiva de técnicas de preprocesamiento de datos (selección de características y balanceo de clases) aplicadas al conjunto de datos BugHunter para la predicción de fallas de software.
    \item Se identifican las métricas de código más relevantes para la predicción de fallas, demostrando que es posible reducir la dimensionalidad de los conjuntos de datos sin afectar la precisión de los modelos.
    \item Se analiza el impacto de diferentes algoritmos de balanceo de clases en la precisión de los modelos de predicción.
    \item Se compara el desempeño de modelos entrenados con conjuntos de datos individuales y conjuntos de datos unificados (proyecto resumen), aportando evidencia sobre la viabilidad de integrar métricas de múltiples proyectos.
    \item Se proporciona una línea base replicando experimentos previos y se proponen mejoras metodológicas que pueden ser adoptadas en investigaciones futuras.
\end{itemize}

El resto del documento está organizado de la siguiente manera: el Capítulo II introduce trabajos relacionados con el preprocesamiento de datos para la predicción de errores de software. El Capítulo III describe el conjunto de datos BugHunter, incluyendo materiales y métodos utilizados. El Capítulo IV analiza resultados y responde a las preguntas de investigación. Finalmente, el Capítulo V presenta las principales conclusiones, identificando sus limitaciones y trabajos futuros.

\newpage